\documentclass[12pt,a4paper]{article}

\usepackage[margin=2.5cm]{geometry}
\usepackage{booktabs}
\usepackage{tabularx}
\usepackage{array}
\usepackage{hyperref}
\usepackage{enumitem}
\usepackage{fancyhdr}
\usepackage{parskip}
\usepackage{titlesec}
\usepackage{xcolor}

\definecolor{accent}{RGB}{30,80,160}

\hypersetup{colorlinks=true, linkcolor=accent, urlcolor=accent}

\titleformat{\section}{\large\bfseries\color{accent}}{}{0em}{}[\titlerule]
\titleformat{\subsection}{\normalsize\bfseries\color{accent}}{}{0em}{}

\pagestyle{fancy}
\fancyhf{}
\lhead{\textbf{WanderSync} \textit{Task 3: Architectural Tactics}}
\rhead{Team 15}
\rfoot{\thepage}
\renewcommand{\headrulewidth}{0.4pt}

\begin{document}

\begin{center}
  {\LARGE\bfseries Task 3: Architectural Tactics}\\[0.4em]
  {\large WanderSync: Dynamic Itinerary, Disruption Management and Social Travel Hub}\\[0.2em]
  {\normalsize Team 15}
\end{center}

\vspace{1em}

% ================================================================
\section{Overview}
% ================================================================

Architectural tactics are design decisions that directly influence the system's response to a
quality-attribute stimulus. Each tactic below is selected to address a specific non-functional
requirement (NFR) documented in Task 1. The five tactics collectively span the performance,
reliability, security, and privacy quality attributes of WanderSync.

% ================================================================
\section{Architectural Tactics}
% ================================================================

% ----------------------------------------------------------------
\subsection{Tactic 1: Caching}
% ----------------------------------------------------------------

\textbf{NFR addressed:} NFR1 -- Performance (itinerary page loads in under 3 seconds; external
API calls complete within 5 seconds).

\textbf{Mechanism in WanderSync:}
The Itinerary Management Service caches fully-assembled itinerary objects in Redis, keyed by
itinerary identifier. On an inbound timeline request the service performs a cache lookup first;
only on a cache miss does it re-query the database and re-aggregate provider data. A separate
short-lived cache layer (60-second TTL) in the Booking Aggregation Service stores raw provider
responses, absorbing burst traffic without redundant external API calls.

The tactic directly satisfies the 3-second render requirement by eliminating database and
provider round-trips on the hot path. Cache entries are invalidated on itinerary mutation (booking
added, disruption applied) to maintain consistency.

\textbf{Trade-off:} Cached data may be transiently stale after a disruption event fires before
invalidation completes. This is accepted because the EDA pipeline delivers \texttt{DisruptionEvent}
messages to the Itinerary Service within 10 seconds, triggering cache invalidation promptly.

% ----------------------------------------------------------------
\subsection{Tactic 2: Retry with Circuit Breaker}
% ----------------------------------------------------------------

\textbf{NFR addressed:} NFR1 -- Performance (external API calls under 5 seconds) and NFR2 --
Reliability (graceful error handling; no crash on external API failure; disruption notification
within 10 seconds).

\textbf{Mechanism in WanderSync:}
Every call from the Booking Aggregation Service to an external provider adapter is wrapped in a
\texttt{CircuitBreakerProxy}. On a transient failure the proxy retries up to two times with
exponential backoff (500\,ms, 1\,s). If the failure threshold is exceeded (three consecutive
failures within a 30-second window), the circuit opens and the aggregation layer falls back to
cached data or marks that provider as temporarily unavailable, returning a partial response with
a user-visible status indicator instead of propagating an exception.

The Disruption Detection Service applies the same retry policy when polling external
flight-tracking feeds, ensuring that transient API unavailability does not silently drop detection
events and degrade accuracy below the 17/20 threshold.

\textbf{Trade-off:} The maximum retry window (up to 3 seconds with backoff) is bounded well within
the 5-second API latency NFR. Open-circuit fallback means a subset of provider data may be absent
from the itinerary view during an outage, which is communicated to the user rather than surfaced
as an error.

% ----------------------------------------------------------------
\subsection{Tactic 3: Authenticate and Authorise at the Gateway}
% ----------------------------------------------------------------

\textbf{NFR addressed:} NFR3 -- Security (100\% of protected endpoints require a valid JWT with
24-hour expiry; all traffic over HTTPS; no sensitive data over plain HTTP).

\textbf{Mechanism in WanderSync:}
The API Gateway intercepts every inbound request and validates the JWT signature and expiry
before forwarding to any downstream service. TLS is terminated at the gateway; backend services
communicate over an internal network without additional TLS overhead. No downstream service
accepts client requests directly -- they receive only internally forwarded requests from the
gateway, which appends a decoded identity header (\texttt{X-User-Id}, \texttt{X-User-Role})
after token validation.

The gateway enforces \textbf{authentication (AuthN)} and \textbf{coarse-grained role-based
authorisation (AuthZ)} — for example, confirming that the caller holds the \texttt{TRAVELLER}
or \texttt{ADMIN} role before allowing access to a protected route. This prevents unauthenticated
access and coarse role violations in one central location.

\textbf{Fine-grained resource authorisation} — such as verifying that the authenticated user
actually owns the itinerary they are attempting to edit — \textbf{must remain in the domain
services}. The API Gateway has no access to the Itinerary or Social databases; it cannot
determine whether User A owns Itinerary B. Delegating resource ownership checks entirely to
the gateway would introduce an Insecure Direct Object Reference (IDOR) vulnerability.
Downstream services verify resource ownership using the forwarded identity header before
executing any mutating operation.

\textbf{Trade-off:} The gateway becomes a single enforcement point for AuthN and coarse AuthZ,
which means a gateway misconfiguration could expose all services. This risk is mitigated by
keeping the auth middleware minimal and independently testable. The ownership-check responsibility
in domain services is explicit and tested at the integration level.

% ----------------------------------------------------------------
\subsection{Tactic 4: Event Queue with Persistence}
% ----------------------------------------------------------------

\textbf{NFR addressed:} NFR2 -- Reliability (disruption notification within 10 seconds of event
arrival; at least 85\% detection accuracy across 20 simulated test cases).

\textbf{Mechanism in WanderSync:}
The Disruption Detection Service publishes \texttt{DisruptionEvent} messages to a Redis Streams
channel rather than a fire-and-forget Pub/Sub topic. Redis Streams persists messages to an
append-only log with consumer group semantics: if the Itinerary Management or Notification
consumer is briefly unavailable, it replays missed events upon reconnection without loss.

The persisted event stream also serves as the ground truth for accuracy measurement: the 20
simulated test cases can be replayed against the log to count correctly identified and notified
disruptions deterministically.

\textbf{Trade-off:} Redis Streams introduces slightly more operational complexity than plain
Pub/Sub (consumer group management, stream trimming). This is accepted because persistence is
non-negotiable for the 85\% accuracy NFR and the 10-second delivery window under consumer
restart scenarios.

% ----------------------------------------------------------------
\subsection{Tactic 5: Geospatial Masking}
% ----------------------------------------------------------------

\textbf{NFR addressed:} NFR4 -- Privacy (zero instances of precise location leakage before
bilateral opt-in; city-level region label only before consent; location data not persisted beyond
the active session unless explicitly saved).

\textbf{Mechanism in WanderSync:}
The Social Synchronisation Service applies a \texttt{LocationMaskingFilter} to every outbound
social query response before serialisation. The filter consults the consent registry for the pair
of users involved in the query. If bilateral consent is absent, the filter truncates precise
coordinates to the city centroid (bounding-box snap to nearest city-level geohash). Raw
coordinates never leave the service boundary without a confirmed bilateral opt-in record.

Session-scoped location data is stored exclusively in memory and purged at session termination
unless the user explicitly invokes the save action, preventing residual location leakage in
persistent storage.

\textbf{Trade-off:} City-level masking reduces the spatial resolution of matching results for
non-consented users, potentially surfacing false positives (users in the same city but not truly
nearby). This is an intentional privacy-utility trade-off: precision is gated behind explicit
consent.

\end{document}
